% Options for packages loaded elsewhere
\PassOptionsToPackage{unicode}{hyperref}
\PassOptionsToPackage{hyphens}{url}
\documentclass[
]{article}
\usepackage{xcolor}
\usepackage[margin=1in]{geometry}
\usepackage{amsmath,amssymb}
\setcounter{secnumdepth}{-\maxdimen} % remove section numbering
\usepackage{iftex}
\ifPDFTeX
  \usepackage[T1]{fontenc}
  \usepackage[utf8]{inputenc}
  \usepackage{textcomp} % provide euro and other symbols
\else % if luatex or xetex
  \usepackage{unicode-math} % this also loads fontspec
  \defaultfontfeatures{Scale=MatchLowercase}
  \defaultfontfeatures[\rmfamily]{Ligatures=TeX,Scale=1}
\fi
\usepackage{lmodern}
\ifPDFTeX\else
  % xetex/luatex font selection
\fi
% Use upquote if available, for straight quotes in verbatim environments
\IfFileExists{upquote.sty}{\usepackage{upquote}}{}
\IfFileExists{microtype.sty}{% use microtype if available
  \usepackage[]{microtype}
  \UseMicrotypeSet[protrusion]{basicmath} % disable protrusion for tt fonts
}{}
\makeatletter
\@ifundefined{KOMAClassName}{% if non-KOMA class
  \IfFileExists{parskip.sty}{%
    \usepackage{parskip}
  }{% else
    \setlength{\parindent}{0pt}
    \setlength{\parskip}{6pt plus 2pt minus 1pt}}
}{% if KOMA class
  \KOMAoptions{parskip=half}}
\makeatother
\usepackage{color}
\usepackage{fancyvrb}
\newcommand{\VerbBar}{|}
\newcommand{\VERB}{\Verb[commandchars=\\\{\}]}
\DefineVerbatimEnvironment{Highlighting}{Verbatim}{commandchars=\\\{\}}
% Add ',fontsize=\small' for more characters per line
\usepackage{framed}
\definecolor{shadecolor}{RGB}{248,248,248}
\newenvironment{Shaded}{\begin{snugshade}}{\end{snugshade}}
\newcommand{\AlertTok}[1]{\textcolor[rgb]{0.94,0.16,0.16}{#1}}
\newcommand{\AnnotationTok}[1]{\textcolor[rgb]{0.56,0.35,0.01}{\textbf{\textit{#1}}}}
\newcommand{\AttributeTok}[1]{\textcolor[rgb]{0.13,0.29,0.53}{#1}}
\newcommand{\BaseNTok}[1]{\textcolor[rgb]{0.00,0.00,0.81}{#1}}
\newcommand{\BuiltInTok}[1]{#1}
\newcommand{\CharTok}[1]{\textcolor[rgb]{0.31,0.60,0.02}{#1}}
\newcommand{\CommentTok}[1]{\textcolor[rgb]{0.56,0.35,0.01}{\textit{#1}}}
\newcommand{\CommentVarTok}[1]{\textcolor[rgb]{0.56,0.35,0.01}{\textbf{\textit{#1}}}}
\newcommand{\ConstantTok}[1]{\textcolor[rgb]{0.56,0.35,0.01}{#1}}
\newcommand{\ControlFlowTok}[1]{\textcolor[rgb]{0.13,0.29,0.53}{\textbf{#1}}}
\newcommand{\DataTypeTok}[1]{\textcolor[rgb]{0.13,0.29,0.53}{#1}}
\newcommand{\DecValTok}[1]{\textcolor[rgb]{0.00,0.00,0.81}{#1}}
\newcommand{\DocumentationTok}[1]{\textcolor[rgb]{0.56,0.35,0.01}{\textbf{\textit{#1}}}}
\newcommand{\ErrorTok}[1]{\textcolor[rgb]{0.64,0.00,0.00}{\textbf{#1}}}
\newcommand{\ExtensionTok}[1]{#1}
\newcommand{\FloatTok}[1]{\textcolor[rgb]{0.00,0.00,0.81}{#1}}
\newcommand{\FunctionTok}[1]{\textcolor[rgb]{0.13,0.29,0.53}{\textbf{#1}}}
\newcommand{\ImportTok}[1]{#1}
\newcommand{\InformationTok}[1]{\textcolor[rgb]{0.56,0.35,0.01}{\textbf{\textit{#1}}}}
\newcommand{\KeywordTok}[1]{\textcolor[rgb]{0.13,0.29,0.53}{\textbf{#1}}}
\newcommand{\NormalTok}[1]{#1}
\newcommand{\OperatorTok}[1]{\textcolor[rgb]{0.81,0.36,0.00}{\textbf{#1}}}
\newcommand{\OtherTok}[1]{\textcolor[rgb]{0.56,0.35,0.01}{#1}}
\newcommand{\PreprocessorTok}[1]{\textcolor[rgb]{0.56,0.35,0.01}{\textit{#1}}}
\newcommand{\RegionMarkerTok}[1]{#1}
\newcommand{\SpecialCharTok}[1]{\textcolor[rgb]{0.81,0.36,0.00}{\textbf{#1}}}
\newcommand{\SpecialStringTok}[1]{\textcolor[rgb]{0.31,0.60,0.02}{#1}}
\newcommand{\StringTok}[1]{\textcolor[rgb]{0.31,0.60,0.02}{#1}}
\newcommand{\VariableTok}[1]{\textcolor[rgb]{0.00,0.00,0.00}{#1}}
\newcommand{\VerbatimStringTok}[1]{\textcolor[rgb]{0.31,0.60,0.02}{#1}}
\newcommand{\WarningTok}[1]{\textcolor[rgb]{0.56,0.35,0.01}{\textbf{\textit{#1}}}}
\usepackage{graphicx}
\makeatletter
\newsavebox\pandoc@box
\newcommand*\pandocbounded[1]{% scales image to fit in text height/width
  \sbox\pandoc@box{#1}%
  \Gscale@div\@tempa{\textheight}{\dimexpr\ht\pandoc@box+\dp\pandoc@box\relax}%
  \Gscale@div\@tempb{\linewidth}{\wd\pandoc@box}%
  \ifdim\@tempb\p@<\@tempa\p@\let\@tempa\@tempb\fi% select the smaller of both
  \ifdim\@tempa\p@<\p@\scalebox{\@tempa}{\usebox\pandoc@box}%
  \else\usebox{\pandoc@box}%
  \fi%
}
% Set default figure placement to htbp
\def\fps@figure{htbp}
\makeatother
\setlength{\emergencystretch}{3em} % prevent overfull lines
\providecommand{\tightlist}{%
  \setlength{\itemsep}{0pt}\setlength{\parskip}{0pt}}
\usepackage{bookmark}
\IfFileExists{xurl.sty}{\usepackage{xurl}}{} % add URL line breaks if available
\urlstyle{same}
\hypersetup{
  pdftitle={Final\_Report},
  pdfauthor={Brisa Banuelos and Quynh Pham},
  hidelinks,
  pdfcreator={LaTeX via pandoc}}

\title{Final\_Report}
\author{Brisa Banuelos and Quynh Pham}
\date{2026-04-19}

\begin{document}
\maketitle

\section{Introduction}\label{introduction}

This report goes through a variety of health variables and shows how
they could be connected. These variables include quantitative variables,
such as, age, weight, and height. Qualitative variables include sex of
sample adult, race/ethnicity, educational level, general health status,
and life satisfaction.

\section{Methods}\label{methods}

After loading the data into R studio, the data was cleaned and checked
for missing values in the variables. After cleaning the data, there were
multiple methods used to display the connection between certain
variables. Some of these methods included showing a summary of the
statistics for the quantitative variables, creating frequency tables for
qualitative variables, creating bar plots, box plots, and scatter plots
to compare the two types of variables.

\section{Results}\label{results}

To begin, we first have to visualize the distribution of the health
variables of the participants. When looking at age distribution through
a histogram, there was a somewhat even distribution of the participant's
ages. The more ``popular'' ages of the participants included around 63,
71, and 55 years old. The weight distribution of the participants showed
a histogram slightly skewed to the right. Some of the ``popular''
weights included around 180, 160, and 149 pounds. The height
distribution appears approximately normal, with no strong skew. A few of
the ``popular'' heights included being around 64, 66, and 67 inches.
Lastly, the data appeared to have slightly more female participants than
male participants. The scatter plots that compare height vs weight show
that as the height of a participant increases, so does the participant's
weight.

\begin{Shaded}
\begin{Highlighting}[]
\FunctionTok{cor}\NormalTok{(nhis}\SpecialCharTok{$}\NormalTok{HEIGHTTC\_A, nhis}\SpecialCharTok{$}\NormalTok{WEIGHTLBTC\_A)}
\end{Highlighting}
\end{Shaded}

\begin{verbatim}
## [1] 0.5023037
\end{verbatim}

The correlation between height and weight is positive, indicating a
linear relationship where taller individuals tend to weigh more. When
looking at the enhanced scatter plot that compares height and weight by
sex and educational level, it shows that as height increases, weight
also increases, and that the majority of male participants are heavier
and taller than the female participants. Additionally, a large
proportion of participants fall into the high school graduate level and
college graduate or better level. Lastly, the correlation matrix showed
a weaker or slightly negative relationship between age and weight, a
negative relationship between age and height, and a positive
relationship between weight and height.

\section{Discussion}\label{discussion}

After reviewing the results from the data collected and visuals created,
there is a clear relationship between weight and height, which is, as a
participant's height increases, their weight increases as well.
Additionally, males were more likely to have a higher weight and height
compared to the female participants, as well a large proportion of
participants either being high school graduates or college graduates.
The findings suggest that demographic factors such as sex and education
level are associated with differences in height and weight, although
further analysis would be needed to determine causation.

\section{Conclusion}\label{conclusion}

By reviewing the data collected, the associations between certain health
variables, such as weight, height, sex, and educational level, are able
to be seen. For example, there was a positive relationship between
height and weight consistently throughout the different visualizations
created, such as scatterplots and the correlation matrix.

\end{document}
